% !PRINTED 459
\weberpage{459}
%<<CONT
\[
24\,P_7,\ 56\,P_3,\ 84\,\bar P_2 \ \text{ on the curve } f.
\]
\[
21\,P_8,\ 42\,P_4,\ 28\,P_6 \ \text{ not on the curve } f.
\]
All other points of the curve $f$ are carried by the substitutions of
the group into 168 distinct points.

From this it follows, moreover, immediately that the curve $f$ can
have no double point, and hence still less can it split into curves of
lower degree. For suppose it had a double point; then all its image
points would also have to be double points, and since a curve of the
fourth order, even when it splits, cannot have more than six double
points unless it has infinitely many double points, i.\,e. doubly
counted parts of the curve, $f$ would have to be the square of a
quadratic form. The root of $f$ would then also have to be an
invariant, whereas there are no quadratic invariants.

One can, incidentally, also easily show directly that the fundamental
curve has no double point; for the conditions for a double point
are:
\[
f'(x_1) = 3x_1^2x_3 + x_2^3 = 0,
\]
\[
f'(x_2) = 3x_2^2x_1 + x_3^3 = 0,
\]
\[
f'(x_3) = 3x_3^2x_2 + x_1^3 = 0,
\]
and these equations cannot be satisfied otherwise than when
$x_1, x_2, x_3$ vanish (\S.\,86).

For the sake of brevity, we shall call a system of points, each of
which can be transformed into every other by substitutions of the
group, a \gesperrt{system of connected points}.

\S.\,123.

\gesperrt{The higher invariants.}

Further invariants of our group can be derived, according to Theorem 4.,
\S.\,40, by forming covariants of the form $f$. We fix our attention
on the covariants discussed in general terms in \S.\,87, and first form
the \gesperrt{Hesse} covariant
\weq{1}{\Delta = \frac{1}{54}
\begin{vmatrix}
f''(x_1,x_1), & f''(x_1,x_2), & f''(x_1,x_3) \\
f''(x_2,x_1), & f''(x_2,x_2), & f''(x_2,x_3) \\
f''(x_3,x_1), & f''(x_3,x_2), & f''(x_3,x_3)
\end{vmatrix},}
%CONT>>
